\section{Conclusion}
\label{sec:conclusion}

Critics of psychiatric deinstitutionalization predicted that closing
hospitals would kill patients. In Brazil, where the Reforma Psiquiátrica
closed specialized psychiatric hospitals across two decades, we find no
evidence that it did. Among the populations that actually depended on
the closing hospitals, suicide mortality changed by $\valSuicWAtt$ per
$100{,}000$ (95\% CI $[\valSuicWLo, \valSuicWHi]$) and self-harm mortality was
likewise null, with parallel pre-trends that hold on data extending back
to $2010$. The interval rules out the large increases the reform's
critics feared; it does not establish a precise zero, and we are careful
not to claim one. The catastrophe did not happen to the people who bore
the closures.

That answer is credible only because we measure exposure correctly.
Identifying who depended on a closing hospital by geographic proximity,
as the closure literature does, misclassifies $\valNmisclass$ of the $\valNunionPairs$
municipality--closure pairs in our sample---missing the populations that
crossed borders for specialized care and flagging neighbors who never
used the facility. A patient-flow exposure rule, built from $\valNadmissions$
million SUS admissions and validated by a sharp post-closure fall in
travel burden, recovers the population a distance-based analysis would
have missed. The measurement correction is the paper's portable
contribution: in any single-payer system with origin--destination claims
data and specialized care concentrated in regional hubs, distance-based
closure studies are estimating effects on the wrong population.

We hold the same evidentiary standard against ourselves. Preventable
hospitalizations rise after closure, but the full pre-periods show they
were already rising before it; we report the trajectory and decline to
read it causally rather than present a pre-existing trend as an effect.
The discipline matters for how the null should be read: it is the
product of an analysis that identified what the data could bear and set
aside what they could not.

The estimates are local to Brazil's psychiatric-reform context and the
SUS-using population, and the closure sample is modest, so we do not
claim that deinstitutionalization is harmless everywhere or that subtle
effects are absent. What we establish is narrower and, for a contested
global policy, useful: when Brazil closed the psychiatric hospitals,
the mortality catastrophe its critics predicted did not appear in the
populations that depended on them.
